% SPDX-License-Identifier: Apache-2.0
\section{The new component}
\label{sec:component}

No existing tool reads a FIDO key through usbfs and relays it. The smallest
new component, \code{ffwd-usbfs}, fills that gap. Table~\ref{tab:component}
gives its design. We wrote it after the probe passed on a real Chromebook
(Section~\ref{sec:testplan}), and it passed the end-to-end test T4.

\begin{table}[t]
\centering
\caption{Design of \code{ffwd-usbfs}}
\label{tab:component}
\scriptsize
\begin{tabular}{@{}>{\raggedright\arraybackslash}p{15mm}>{\raggedright\arraybackslash}p{63mm}@{}}
\toprule
Item & Design \\
\midrule
Purpose & Key side of fido-forward for a host with a usbfs node and no \code{hidraw} node \\
Language & Go 1.27, \code{CGO\_ENABLED=0}, built with Bazel (\code{rules\_go}) \\
Libraries & \code{golang.org/x/sys/unix}, and the packages \code{ctaphid} and \code{usbfs} of this study \\
Command & \code{ffwd-usbfs [-d NODE] [--allow-management] HOST} \\
Transport & Starts \code{ssh -T HOST fido-forward --sink} with pipes, as fido-forward does; under \code{sudo} it gives up root before \code{ssh} starts \\
Wire & fido-forward v1: read \code{FFWD\textbackslash x01}, write it back, then raw 64-byte reports \\
USB & Find the HID interface with page 0xF1D0; detach \code{usbhid}; claim it; one reader of interrupt IN, one writer of interrupt OUT \\
Policy & Judge each request by its first packet; answer a refused CTAP2 command with status \code{0x27} (operation denied) and a refused vendor command with \code{CTAPHID\_ERROR}; drop the rest of the message; never send it to the key \\
Teardown & \code{CTAPHID\_CANCEL} on each channel, release, attach \code{usbhid} again \\
Size & 586 lines, and 214 lines of tests \\
\bottomrule
\end{tabular}
\end{table}

\subsection{Parts that exist now}

The study's workspace already holds the parts that need the most care. Both
are pure Go with unit tests, and both build with Bazel.

\paragraph{\code{ctaphid}} It splits a message into 64-byte reports and joins
them again. Its \code{Policy} type decides which CTAP2 commands pass.
Listing~\ref{lst:policy} shows the decision.

\begin{lstlisting}[style=go,caption={The command policy in package \code{ctaphid}},label={lst:policy},float=tb]
// Allowed: does the policy pass m to the key?
func (p Policy) Allowed(m Message) bool {
  if m.Command != CmdCBOR || len(m.Payload) == 0 {
    return true // framing (INIT, PING, CANCEL) or U2F
  }
  switch m.Payload[0] {
  case AuthMakeCredential, AuthGetAssertion, AuthGetInfo,
    AuthClientPIN, AuthGetNextAssertion, AuthSelection:
    return true
  case AuthReset, AuthCredentialManagement, AuthBioEnrollment,
    AuthConfig, AuthLargeBlobs, AuthCredMgmtPreview,
    AuthBioEnrollPreview:
    return p.AllowManagement
  default:
    return false // unknown commands are refused
  }
}
\end{lstlisting}

\paragraph{\code{usbfs}} It opens a usbfs node, parses the descriptors,
reads a HID report descriptor, detaches and attaches the kernel driver,
claims an interface, and runs interrupt transfers. A unit test checks that
each \code{ioctl} number that the code computes from its structure sizes
equals the number in \code{<linux/usbdevice\_fs.h>}.

\paragraph{\code{fido-probe}} It uses both packages. It finds the YubiKey,
detaches \code{usbhid} from the FIDO interface, sends \code{CTAPHID\_INIT}
and \code{getInfo}, and attaches the driver again. Both commands read only.

\paragraph{\code{relay}} It is the key side of fido-forward's wire
protocol. It has no USB code, so unit tests run it against a fake key. They
check that a reset of three packets never reaches the key, that the far side
gets status \code{0x27}, and that every used channel gets
\code{CTAPHID\_CANCEL} at the end.

\begin{tip}
The first version of \code{ffwd-usbfs} reads the interrupt IN endpoint in
a loop with a 200\,ms timeout. A later version can keep an IN transfer
pending at all times with asynchronous URBs (\code{USBDEVFS\_SUBMITURB} and
\code{USBDEVFS\_REAPURB}).
\end{tip}
